\section{第二机会曲线的三道门槛}

已启动方向进入核心池必须同时满足三项要求：第一，20日或60日价格趋势已经转强，且有资金或成交确认；第二，增长叙事能连接到2026--2028收入、净利润和EPS，而不是只有题材空间；第三，按7月10日现价重算后，AStock多锚目标仍至少高出约20\%。工业富联在H1预告后同时满足三项要求，新增为第四只B篮子核心。只满足前两项、但价格已充分反映的先进封装，不进入核心；只满足趋势、但利润桥不完整的传媒/AI应用，也不进入核心。

\begin{exhibitbox}[Exhibit 9：已启动高空间核心与条件观察]
\begin{tabularx}{\textwidth}{L{2.3cm}L{2.3cm}R{1.6cm}R{1.7cm}R{1.7cm}X}
\toprule
\textbf{方向} & \textbf{核心/代表} & \textbf{20日涨跌} & \textbf{现价} & \textbf{目标空间} & \textbf{结论} \\
\midrule
AI网络设备 & 中兴通讯 & +7.2\% & 40.53 & +20.3\% & 已启动核心，等H2利润和现金流验证 \\
AI基础设施 & 工业富联 & -4.7\% & 66.27 & +24.2\% & H1预告与预告后研报闭环，等H2平台量产 \\
AI存储 & 江波龙 & +14.1\% & 587.60 & +36.5\% & 盈利弹性最大，只做回撤配置 \\
创新药 & 恒瑞医药 & +18.7\% & 55.75 & +34.4\% & 增长质量最好，SOTP仍有空间 \\
国防军工 & 中航光电等 & 板块启动 & --- & 未过核心门槛 & Q1利润承压，等待中报与订单恢复 \\
先进封装 & 华天/长电 & \shortstack{+51\%\\+41\%} & \shortstack{25.31\\101.11} & 空间不足 & 产业正确、价格提前反映 \\
传媒/AI应用 & 板块 & 周涨2.01\% & --- & 无完整公司桥 & 主题观察，不给核心估值信用 \\
\bottomrule
\end{tabularx}
\end{exhibitbox}

\section{AI网络设备：中兴通讯从连接设备向算力底座重估}

中兴通讯在7月6日至10日一周获得约49亿元主力净流入，7月9日放量上涨后，7月10日以186亿元成交额高位震荡，启动信号明确。基本面并非没有瑕疵：2026Q1收入349.9亿元、同比+6.1\%，归母净利润13.1亿元、同比下降46.6\%，经营现金流约-19.8亿元。但算力产品收入占比已经提升至27\%，说明第二增长曲线正在形成。

本报告不直接采用旧券商92.6亿元净利预测，而使用当前公开一致预期约71亿元、EPS 1.48元。现价40.53元对应约27.4倍2026E PE，低于AStock基准33倍。最终目标48.75元不是单纯估值拔高：升级条件是算力收入占比保持25\%以上、H2利润恢复正增长、经营现金流改善；否则只保留算力期权，不给基准倍数。

\section{AI基础设施：工业富联从需求锚升级为正式核心}

工业富联H1预告归母净利234--244亿元，中值239亿元，同比约+97\%；隐含Q2约133.05亿元，较Q1增长25.6\%。公司同时披露云服务商AI服务器收入同比增长超过230\%、800G以上数据中心交换机出货同比增长1.4倍，使收入、产品结构和利润形成直接证据链。

华泰预告后维持2026E净利560.1亿元、EPS 2.82元和93元目标；华创给出更高的645.97亿元净利，但未在归档页面披露目标价。AStock以华泰为基准分母、华创为牛市交叉验证，最终目标82.29元，较66.27元现价空间24.2\%。动作是业绩兑现后的回撤配置，而非把工业富联误判为低位标的。

\section{AI存储：江波龙已经启动，但利润分母上修更快}

江波龙是已启动篮子中弹性最大、争议也最大的标的。股价60日上涨约66\%，7月10日回落至587.60元；与此同时，官方H1预告给出收入220--250亿元、归母净利润92--110亿元，Q2隐含利润继续环比提升。市场价格上涨很快，但盈利分母上修更快。

按2026E净利润200亿元、EPS 47.28元，现价约12.4倍PE。AStock基准价值900元，最终多锚目标802.30元。这个空间不能被理解为“周期消失”：Q1经营现金流为负、库存绝对规模大，摩根士丹利虽然把目标价上调至673元，仍维持中性，核心担忧正是模组厂供应分配、库存红利消耗和消费端价格见顶。因此动作定义为\textbf{高弹性回撤配置}，而不是突破追涨。

\section{创新药：恒瑞医药的启动来自收入结构重构}

恒瑞医药2026Q1收入81.41亿元、归母净利润22.82亿元；创新药销售45.26亿元、同比+25.75\%，占药品销售收入61.69\%，另确认BD收入7.87亿元。华泰预计2026--2028归母净利润94/117.9/149.8亿元，并给出87.14元SOTP目标；高盛目标79.20元。

与纯主题创新药不同，恒瑞的增长桥已经覆盖产品销售、BD、毛利率和费用率。AStock使用78元基准价值和74.91元多锚目标，相对现价空间约34.4\%。最大风险不是“有没有创新药”，而是创新药全年30\%增长目标、BD确认节奏和关键临床里程碑能否持续兑现。

\section{为什么军工、先进封装和传媒只做条件观察}

\begin{itemize}
  \item \textbf{军工}：行业资金和事件启动明确，但中航光电2026Q1归母净利润同比下降37.8\%、经营现金流为负。若中报军品订单与利润恢复，才可从板块确认升级为核心公司估值。
  \item \textbf{先进封装}：华天科技和长电科技20日涨幅约51\%和41\%，分别处于一年价格高位。产业空间很大，但短期价格空间和安全边际不足。
  \item \textbf{传媒/AI应用}：板块资金改善，但缺少能把AI应用用户/付费/版权价值连接到公司EPS的完整公开证据，暂不给投资级增长信用。
\end{itemize}
